\input{preamble.tex}

\title{Лабораторная работа № 3. \\ Планирование локальной сети организации}
\author{Сущенко Алина \\ НПИбд-01-23}

\institute{Российский университет дружбы народов имени Патриса Лумумбы}
\date{2026}

\begin{document}

\frame{\titlepage}

\begin{frame}
\frametitle{Цель работы}
\begin{itemize}
    \item Ознакомление с принципами планирования локальной сети организации на примере учебного заведения с двумя территориально распределенными площадками.
    \item Получение навыков разработки иерархической структуры сети (уровни ядра, распределения и доступа).
    \item Освоение методик планирования VLAN, IP-адресации и портов подключения оборудования.
\end{itemize}
\end{frame}

\begin{frame}
\frametitle{Физическая структура сети (L1)}
\begin{itemize}
    \item Разработана физическая схема сети в соответствии с иерархической моделью.
    \item На схеме указаны все устройства и номера портов подключения.
    \item \textbf{Уровень ядра:} маршрутизатор Cisco 2811 (msk-ansusenko-gw-1)
    \item \textbf{Уровень распределения:} коммутаторы Cisco 2960 (msk-ansusenko-sw-1, sw-2, sw-3, sw-4)
    \item \textbf{Уровень доступа:} коммутатор Cisco 2950 (msk-ansusenko-sw-1)
\end{itemize}
\begin{center}
    \includegraphics[width=0.7\textwidth]{../images/01.png}
\end{center}
\caption{Физические устройства сети с номерами портов (L1)}
\end{frame}

\begin{frame}
\frametitle{Схема VLAN (L2)}
\begin{itemize}
    \item Сеть разделена на логические сегменты для изоляции трафика различных групп пользователей.
    \item Магистральные каналы (trunk) передают трафик нескольких VLAN между коммутаторами.
    \item Порты доступа (access) привязаны к конкретным VLAN для подключения оконечных устройств.
\end{itemize}
\begin{table}[h]
\centering
\caption{План VLAN}
\begin{tabular}{|c|c|c|}
\hline
\textbf{№ VLAN} & \textbf{Имя VLAN} & \textbf{Примечание} \\
\hline
1 & default & Не используется \\
2 & management & Управление устройствами \\
3 & servers & Серверная ферма \\
101 & dk & Дисплейные классы (ДК) \\
102 & departments & Кафедры \\
103 & adm & Администрация \\
104 & other & Другие пользователи \\
\hline
\end{tabular}
\end{table}
\end{frame}

\begin{frame}
\frametitle{Схема VLAN (L2)}
\begin{center}
    \includegraphics[width=0.6\textwidth]{../images/02.png}
\end{center}
\caption{Схема VLAN сети (L2)}
\end{frame}

\begin{frame}
\frametitle{Схема маршрутизации для сети 10.128.0.0/16 (L3)}
\begin{itemize}
    \item Для исходного проекта используется адресное пространство 10.128.0.0/16.
    \item Каждому VLAN выделена отдельная подсеть /24.
    \item Над каждой группой пользователей указаны сеть и адрес шлюза.
\end{itemize}
\begin{center}
    \includegraphics[width=0.8\textwidth]{../images/03.png}
\end{center}
\caption{Схема маршрутизации сети (L3)}
\end{frame}

\begin{frame}
\frametitle{План IP-адресации для сети 10.128.0.0/16}
\begin{table}[h]
\centering
\caption{IP-адресация для 10.128.0.0/16}
\begin{tabular}{|c|c|c|c|}
\hline
\textbf{Подсеть} & \textbf{IP-сеть} & \textbf{VLAN} & \textbf{Шлюз} \\
\hline
Серверная ферма & 10.128.0.0/24 & 3 & 10.128.0.1 \\
Управление & 10.128.1.0/24 & 2 & 10.128.1.1 \\
Point-to-Point & 10.128.2.0/24 & - & 10.128.2.1 \\
ДК & 10.128.3.0/24 & 101 & 10.128.3.1 \\
Кафедры & 10.128.4.0/24 & 102 & 10.128.4.1 \\
Администрация & 10.128.5.0/24 & 103 & 10.128.5.1 \\
Другие & 10.128.6.0/24 & 104 & 10.128.6.1 \\
\hline
\end{tabular}
\end{table}
\end{frame}

\begin{frame}
\frametitle{План IP-адресации для сети 172.16.0.0/12}
\begin{itemize}
    \item Для диапазона 172.16.0.0/12 выбрана базовая сеть 172.16.0.0/16 для сохранения структуры /24 подсетей.
    \item Физическая схема (L1) и схема VLAN (L2) полностью аналогичны предыдущему проекту.
\end{itemize}
\begin{table}[h]
\centering
\caption{IP-адресация для 172.16.0.0/12}
\begin{tabular}{|c|c|c|c|}
\hline
\textbf{Подсеть} & \textbf{IP-сеть} & \textbf{VLAN} & \textbf{Шлюз} \\
\hline
Серверная ферма & 172.16.0.0/24 & 3 & 172.16.0.1 \\
Управление & 172.16.1.0/24 & 2 & 172.16.1.1 \\
Point-to-Point & 172.16.2.0/24 & - & 172.16.2.1 \\
ДК & 172.16.3.0/24 & 101 & 172.16.3.1 \\
Кафедры & 172.16.4.0/24 & 102 & 172.16.4.1 \\
Администрация & 172.16.5.0/24 & 103 & 172.16.5.1 \\
Другие & 172.16.6.0/24 & 104 & 172.16.6.1 \\
\hline
\end{tabular}
\end{table}
\end{frame}

\begin{frame}
\frametitle{План IP-адресации для сети 192.168.0.0/16}
\begin{itemize}
    \item Диапазон 192.168.0.0/16 идеально подходит под структуру /16.
    \item Сохранена та же логическая структура VLAN.
\end{itemize}
\begin{table}[h]
\centering
\caption{IP-адресация для 192.168.0.0/16}
\begin{tabular}{|c|c|c|c|}
\hline
\textbf{Подсеть} & \textbf{IP-сеть} & \textbf{VLAN} & \textbf{Шлюз} \\
\hline
Серверная ферма & 192.168.0.0/24 & 3 & 192.168.0.1 \\
Управление & 192.168.1.0/24 & 2 & 192.168.1.1 \\
Point-to-Point & 192.168.2.0/24 & - & 192.168.2.1 \\
ДК & 192.168.3.0/24 & 101 & 192.168.3.1 \\
Кафедры & 192.168.4.0/24 & 102 & 192.168.4.1 \\
Администрация & 192.168.5.0/24 & 103 & 192.168.5.1 \\
Другие & 192.168.6.0/24 & 104 & 192.168.6.1 \\
\hline
\end{tabular}
\end{table}
\end{frame}

\begin{frame}
\frametitle{План подключения портов оборудования}
\begin{table}[h]
\centering
\caption{План портов (фрагмент)}
\begin{tabular}{|c|c|c|c|}
\hline
\textbf{Устройство} & \textbf{Порт} & \textbf{Назначение} & \textbf{VLAN} \\
\hline
msk-donskaya-gw-1 & f0/0 & UpLink & Trunk 2,3,101-104 \\
msk-donskaya-sw-1 & f0/24 & к gw-1 & Trunk 2,3,101-104 \\
msk-donskaya-sw-1 & g0/1 & к sw-2 & Trunk 2,3 \\
msk-donskaya-sw-1 & f0/1 & к pavlovskaya-sw-1 & Trunk 2,101,104 \\
msk-donskaya-sw-3 & f0/1 & Web-сервер & Access 3 \\
msk-donskaya-sw-4 & f0/1-5 & Пользователи ДК & Access 101 \\
msk-pavlovskaya-sw-1 & f0/1-15 & Пользователи ДК & Access 101 \\
\hline
\end{tabular}
\end{table}
\begin{itemize}
    \item План портов является единым для всех трех проектов, так как физическая топология не меняется.
    \item Четко разделены access-порты (для конечных устройств) и trunk-порты (межкоммутаторные соединения).
\end{itemize}
\end{frame}

\begin{frame}
\frametitle{Контрольные вопросы}
\begin{block}{Вопрос 1: Что такое модель OSI? Какие уровни в ней есть?}
\textbf{OSI (Open Systems Interconnection)} - эталонная модель взаимодействия открытых систем, разработанная ISO для стандартизации сетевых протоколов.
\vspace{0.3cm}

\textbf{Уровни модели OSI (сверху вниз):}
\begin{enumerate}
    \item \textbf{Прикладной (Application)} - взаимодействие пользовательских приложений с сетью (HTTP, FTP, SMTP)
    \item \textbf{Представительский (Presentation)} - преобразование данных, шифрование, сжатие
    \item \textbf{Сеансовый (Session)} - управление сеансами связи, установка/завершение соединений
    \item \textbf{Транспортный (Transport)} - надежная доставка данных, контроль потока (TCP, UDP)
    \item \textbf{Сетевой (Network)} - маршрутизация, логическая адресация (IP, ICMP, маршрутизаторы)
    \item \textbf{Канальный (Data Link)} - физическая адресация (MAC), коммутация (Ethernet, коммутаторы)
    \item \textbf{Физический (Physical)} - передача битов по среде, электрические/оптические сигналы
\end{enumerate}
\end{block}
\end{frame}

\begin{frame}
\frametitle{Контрольные вопросы}
\begin{block}{Вопрос 2: Какие функции выполняет коммутатор?}
\textbf{Коммутатор (Switch)} - устройство канального уровня (L2), выполняющее следующие функции:
\vspace{0.2cm}
\begin{itemize}
    \item \textbf{Коммутация кадров} - передача Ethernet-кадров между портами на основе MAC-адресов
    \item \textbf{Построение таблицы MAC-адресов} - изучение и хранение соответствий MAC-адресов портам
    \item \textbf{Сегментация сети} - разделение сети на коллизионные домены (каждый порт — отдельный домен)
    \item \textbf{Фильтрация трафика} - передача кадров только целевому устройству, а не всем портам
    \item \textbf{Поддержка VLAN} - логическое разделение сети на изолированные широковещательные домены
    \item \textbf{Агрегирование каналов} - объединение нескольких портов для увеличения пропускной способности
    \item \textbf{Обнаружение петель} - поддержка протокола Spanning Tree (STP) для предотвращения широковещательных штормов
\end{itemize}
\end{block}
\end{frame}

\begin{frame}
\frametitle{Контрольные вопросы}
\begin{block}{Вопрос 3: Какие функции выполняет маршрутизатор?}
\textbf{Маршрутизатор (Router)} - устройство сетевого уровня (L3), выполняющее следующие функции:
\vspace{0.2cm}
\begin{itemize}
    \item \textbf{Маршрутизация} - определение оптимального пути передачи пакетов между различными сетями
    \item \textbf{Объединение сетей} - соединение различных IP-подсетей и сетевых технологий
    \item \textbf{NAT (Network Address Translation)} - преобразование частных IP-адресов в публичные
    \item \textbf{Межсетевое экранирование} - фильтрация трафика на основе ACL (Access Control Lists)
    \item \textbf{Поддержка протоколов маршрутизации} - RIP, OSPF, EIGRP, BGP для динамического обновления таблиц
    \item \textbf{Фрагментация пакетов} - разбиение пакетов при передаче между сетями с разным MTU
    \item \textbf{Обеспечение QoS} - приоритезация трафика для критичных приложений
    \item \textbf{Поддержка VPN} - создание защищенных туннелей через публичные сети
\end{itemize}
\end{block}
\end{frame}

\begin{frame}
\frametitle{Контрольные вопросы}
\begin{block}{Вопрос 4: В чём отличие коммутаторов третьего уровня от коммутаторов второго уровня?}
\begin{table}[h]
\centering
\begin{tabular}{|p{4cm}|p{4cm}|}
\hline
\textbf{Коммутатор L2 (уровень 2)} & \textbf{Коммутатор L3 (уровень 3)} \\
\hline
Работает на канальном уровне & Работает на канальном и сетевом уровнях \\
\hline
Принимает решения на основе MAC-адресов & Принимает решения на основе MAC и IP-адресов \\
\hline
Не поддерживает маршрутизацию между VLAN & Поддерживает маршрутизацию между VLAN (межVLAN-роутинг) \\
\hline
Не поддерживает статические/динамические протоколы маршрутизации & Поддерживает протоколы маршрутизации (OSPF, EIGRP, RIP) \\
\hline
Ограничен одним широковещательным доменом (без VLAN) & Может выступать как шлюз по умолчанию для устройств \\
\hline
Более низкая стоимость & Более высокая стоимость \\
\hline
Используется на уровне доступа & Используется на уровне распределения и ядра \\
\hline
\end{tabular}
\end{table}
\end{block}
\end{frame}

\begin{frame}
\frametitle{Контрольные вопросы}
\begin{block}{Вопрос 5: Что такое сетевой интерфейс?}
\textbf{Сетевой интерфейс} - точка соединения устройства с сетью, представленная программными и/или аппаратными средствами.
\vspace{0.2cm}

\textbf{Основные характеристики:}
\begin{itemize}
    \item Имеет уникальный идентификатор на уровне L2 (MAC-адрес)
    \item Может иметь один или несколько IP-адресов на уровне L3
    \item Может быть физическим (сетевая карта) или виртуальным (loopback, VLAN interface, tunnel)
    \item Обеспечивает передачу и прием пакетов данных
\end{itemize}

\textbf{Примеры:}
\begin{itemize}
    \item Физические: Ethernet-порт (eth0, FastEthernet0/0), Wi-Fi адаптер (wlan0)
    \item Виртуальные: Loopback (lo), VLAN-интерфейс (vlan10), интерфейс туннеля (tun0)
    \item На маршрутизаторах: интерфейсы SVI (Switch Virtual Interface) для маршрутизации между VLAN
\end{itemize}
\end{block}
\end{frame}

\begin{frame}
\frametitle{Контрольные вопросы}
\begin{block}{Вопрос 6: Что такое сетевой порт?}
\textbf{Сетевой порт} - термин, имеющий два основных значения в компьютерных сетях:
\vspace{0.2cm}

\textbf{1. Физический порт:}
\begin{itemize}
    \item Аппаратный разъем на сетевом устройстве для подключения кабеля
    \item Примеры: RJ-45 (Ethernet), SFP (оптика), Console-порт, USB
    \item Идентифицируется номером (FastEthernet 0/1, GigabitEthernet 1/0/24)
\end{itemize}

\textbf{2. Программный порт (транспортного уровня):}
\begin{itemize}
    \item Логическая единица для идентификации приложений на устройстве
    \item 16-битное число (от 0 до 65535) в заголовках TCP/UDP
    \item Используется для мультиплексирования соединений
    \item Примеры: порт 80 (HTTP), порт 443 (HTTPS), порт 22 (SSH)
\end{itemize}

\textbf{Важно различать:} физический порт (L1) и программный порт (L4)!
\end{block}
\end{frame}

\begin{frame}
\frametitle{Контрольные вопросы}
\begin{block}{Вопрос 7: Кратко охарактеризуйте технологии Ethernet, Fast Ethernet, Gigabit Ethernet.}
\begin{table}[h]
\centering
\begin{tabular}{|p{2cm}|p{3cm}|p{3cm}|p{1.5cm}|}
\hline
\textbf{Технология} & \textbf{Стандарт} & \textbf{Скорость} & \textbf{Кабель} \\
\hline
\textbf{Ethernet} & 10BASE-T (IEEE 802.3) & 10 Мбит/с & CAT3, CAT5 (до 100 м) \\
\hline
\textbf{Fast Ethernet} & 100BASE-TX (IEEE 802.3u) & 100 Мбит/с & CAT5 (до 100 м) \\
\hline
\textbf{Gigabit Ethernet} & 1000BASE-T (IEEE 802.3ab) & 1000 Мбит/с (1 Гбит/с) & CAT5e, CAT6 (до 100 м) \\
 & 1000BASE-SX (оптика) & 1 Гбит/с & Многомодовое оптоволокно (до 550 м) \\
\hline
\textbf{10 Gigabit Ethernet} & 10GBASE-T (IEEE 802.3an) & 10 Гбит/с & CAT6a, CAT7 (до 100 м) \\
\hline
\end{tabular}
\end{table}

\textbf{Основные различия:}
\begin{itemize}
    \item Скорость передачи данных (от 10 Мбит/с до 10+ Гбит/с)
    \item Методы кодирования сигнала
    \item Типы используемых кабелей и максимальные расстояния
    \item Все технологии сохраняют совместимость формата кадров (frame format)
\end{itemize}
\end{block}
\end{frame}

\begin{frame}
\frametitle{Контрольные вопросы}
\begin{block}{Вопрос 8: Что такое IP-адрес? Определите понятия сеть, подсеть, маска подсети.}
\textbf{IP-адрес (IPv4)} - 32-битное число, уникально идентифицирующее устройство в сети. Записывается в виде четырех десятичных чисел, разделенных точками (например, 192.168.1.1).
\vspace{0.2cm}

\textbf{Основные понятия:}
\begin{itemize}
    \item \textbf{Сеть} - группа устройств, которые могут взаимодействовать напрямую без маршрутизации. Определяется IP-адресом и маской.
    \item \textbf{Подсеть} - логическое разделение крупной сети на более мелкие сегменты для оптимизации маршрутизации и управления.
    \item \textbf{Маска подсети} - 32-битное число, определяющее, какая часть IP-адреса относится к адресу сети, а какая — к адресу узла (host).
\end{itemize}

\textbf{Служебные IP-адреса:}
\begin{itemize}
    \item \textbf{Адрес сети} - все биты хостовой части равны 0 (например, 192.168.1.0/24)
    \item \textbf{Широковещательный адрес (broadcast)} - все биты хостовой части равны 1 (например, 192.168.1.255/24)
    \item \textbf{Адрес шлюза по умолчанию} - обычно первый или последний доступный адрес в подсети
    \item \textbf{Loopback} - 127.0.0.0/8 (обычно 127.0.0.1) для самопроверки
\end{itemize}
\end{block}
\end{frame}

\begin{frame}
\frametitle{Контрольные вопросы}
\begin{block}{Вопрос 8 (продолжение): Пример разбиения сети на подсети}
\textbf{Исходная сеть:} 192.168.1.0/24 (256 адресов, 254 хоста)

\textbf{Задача:} разделить на 4 подсети

\textbf{Решение:} увеличиваем маску на 2 бита (/26)

\begin{table}[h]
\centering
\begin{tabular}{|c|c|c|c|c|}
\hline
\textbf{Подсеть} & \textbf{Адрес сети} & \textbf{Маска} & \textbf{Широковещ.} & \textbf{Хосты} \\
\hline
1 & 192.168.1.0 & 255.255.255.192 & 192.168.1.63 & 62 (1-62) \\
2 & 192.168.1.64 & 255.255.255.192 & 192.168.1.127 & 62 (65-126) \\
3 & 192.168.1.128 & 255.255.255.192 & 192.168.1.191 & 62 (129-190) \\
4 & 192.168.1.192 & 255.255.255.192 & 192.168.1.255 & 62 (193-254) \\
\hline
\end{tabular}
\end{table}

\textbf{Формула расчета:} количество хостов = 2\^(32 - маска) - 2
\begin{itemize}
    \item Для /26: 2\^(32-26) - 2 = 2\^6 - 2 = 64 - 2 = 62 хоста
    \item Вычитаем 2 адреса: адрес сети и широковещательный
\end{itemize}
\end{block}
\end{frame}

\begin{frame}
\frametitle{Контрольные вопросы}
\begin{block}{Вопрос 9: Дайте определение понятию VLAN. Для чего применяется VLAN?}
\textbf{VLAN (Virtual Local Area Network)} - технология, позволяющая логически разделить одну физическую сеть на несколько изолированных широковещательных доменов.
\vspace{0.2cm}

\textbf{Цели применения VLAN:}
\begin{itemize}
    \item \textbf{Безопасность} - изоляция трафика разных отделов (бухгалтерия, HR, администрация)
    \item \textbf{Уменьшение широковещательного трафика} - ограничение broadcast-домена
    \item \textbf{Гибкость} - группировка пользователей по функциям, а не по физическому расположению
    \item \textbf{Упрощение управления} - централизованное применение политик
\end{itemize}

\textbf{Преимущества:}
\begin{itemize}
    \item Логическая группировка пользователей независимо от местоположения
    \item Повышение безопасности за счет изоляции чувствительных данных
    \item Снижение нагрузки на сеть
    \item Простота добавления, перемещения и изменения сотрудников
\end{itemize}
\end{block}
\end{frame}

\begin{frame}
\frametitle{Контрольные вопросы}
\begin{block}{Вопрос 9 (продолжение): Примеры применения VLAN}
\textbf{Пример 1: Многоэтажное офисное здание}
\begin{itemize}
    \item На каждом этаже физически находятся сотрудники разных отделов
    \item VLAN объединяет всех сотрудников отдела продаж, независимо от этажа
    \item Трафик отдела разработки изолирован от трафика бухгалтерии
\end{itemize}

\textbf{Пример 2: Гостевая сеть}
\begin{itemize}
    \item Выделенный VLAN для гостевого Wi-Fi
    \item Доступ только в интернет, без доступа к корпоративным ресурсам
    \item Ограничение скорости и применение дополнительных политик безопасности
\end{itemize}

\textbf{Пример 3: IP-телефония}
\begin{itemize}
    \item Отдельный VLAN для VoIP-трафика
    \item Приоритезация голосового трафика (QoS)
    \item Гарантированное качество связи независимо от загрузки сети данными
\end{itemize}
\end{block}
\end{frame}

\begin{frame}
\frametitle{Контрольные вопросы}
\begin{block}{Вопрос 10: В чём отличие Trunk Port от Access Port?}
\begin{table}[h]
\centering
\begin{tabular}{|p{4cm}|p{4cm}|}
\hline
\textbf{Access Port} & \textbf{Trunk Port} \\
\hline
Принадлежит только одному VLAN & Может передавать трафик нескольких VLAN \\
\hline
Удаляет тег VLAN при отправке конечному устройству & Сохраняет теги VLAN (802.1Q) для идентификации принадлежности \\
\hline
Используется для подключения конечных устройств (ПК, принтеры, IP-телефоны) & Используется для соединения коммутаторов между собой или с маршрутизаторами \\
\hline
Устройство на другом конце обычно не знает о VLAN & Устройство на другом конце должно понимать тегирование 802.1Q \\
\hline
Конфигурация: \texttt{switchport mode access} \newline \texttt{switchport access vlan 10} & Конфигурация: \texttt{switchport mode trunk} \newline \texttt{switchport trunk allowed vlan 10,20,30} \\
\hline
Трафик поступает на устройство без метки VLAN & Трафик передается с меткой VLAN (тегом) \\
\hline
\end{tabular}
\end{table}
\end{block}
\end{frame}

\begin{frame}
\frametitle{Выводы}
\begin{itemize}
    \item Разработана физическая схема сети (L1) для организации с двумя площадками, с четким разделением оборудования по уровням иерархической модели.
    \item Составлен план VLAN (L2), обеспечивающий логическую сегментацию сети и повышение безопасности.
    \item Разработаны планы IP-адресации (L3) для трех частных диапазонов: 10.128.0.0/16, 172.16.0.0/12 и 192.168.0.0/16.
    \item Продемонстрирована универсальность разработанной логической структуры: при смене адресного пространства физическая топология и план VLAN остаются неизменными.
    \item Составлен детальный план подключения интерфейсов оборудования, позволяющий однозначно определить конфигурацию каждого порта.
    \item Полученные навыки могут быть применены при реальном проектировании и развертывании сети предприятия.
\end{itemize}
\end{frame}

\end{document}
